\chapter{Introducción}
\label{cap:introduccion}

\section{Motivación}
\label{sec:intro-motivacion}

El análisis de datos es una disciplina de la informática que me apasiona, así como
el deporte. El análisis de datos aplicado al deporte ha experimentado un crecimiento notable
en la última década. Clubes, televisión, casas de apuestas y compañías de análisis
deportivo invierten cada vez más recursos en extraer información objetiva del
juego a partir del vídeo: distancias recorridas, mapas de calor,
patrones tácticos o detección automática de eventos. Gran parte de esta
información depende de una tarea previa fundamental: saber, en cada instante,
dónde está cada jugador y quién es. Partiendo de esa base, posteriormente se
pueden calcular distintas métricas de interés para el análisis del partido.

Ese problema, conocido en visión por computador como \emph{seguimiento
multiobjeto} o \emph{Multi-Object Tracking} (MOT), es especialmente exigente
en un partido de fútbol: veintidós jugadores más árbitros y balón se mueven de
forma simultánea, se cruzan constantemente, se ocluyen entre sí y comparten
una apariencia muy similar entre compañeros de equipo (misma camiseta, tallas
y complexiones parecidas). Estas condiciones convierten al fútbol en un
escenario particularmente duro para evaluar algoritmos de tracking, mucho más
que los entornos de vigilancia urbana o conducción autónoma para los que
originalmente se diseñaron muchos de estos algoritmos.

Este proyecto nace de aplicar y comparar de forma rigurosa
algoritmos de seguimiento multiobjeto sobre vídeo real de partidos de fútbol
profesional, y comprobar hasta qué punto la información de trayectorias que
producen puede convertirse en metadatos deportivos útiles, sin necesidad de instrumentación adicional
en el terreno de juego, únicamente a partir del vídeo.

\section{Contexto}
\label{sec:intro-contexto}

SoccerNet es un conjunto de datos y una serie de competiciones académicas
centradas en el análisis automático de vídeo de fútbol, con múltiples tareas
publicadas (reconocimiento de acciones, segmentación de cámara, reconocimiento
de dorsales, seguimiento de jugadores, entre otras). Este proyecto se apoya en
la tarea \textbf{SoccerNet-Tracking}, que proporciona secuencias de vídeo de
partidos profesionales descompuestas en frames individuales, junto con
detecciones de personas y balón ya calculadas (\texttt{det.txt}) y, para los
conjuntos de entrenamiento y test, la posición e identidad real de cada
jugador en cada frame (\texttt{gt.txt}), que sirve como referencia para medir
el rendimiento de los algoritmos.

El uso de un dataset público y con un formato estándar permite que los resultados obtenidos en
este proyecto sean comparables con los de otros equipos de investigación, y
que la evaluación se realice con las mismas métricas que utiliza la propia
comunidad de SoccerNet en su tabla de clasificación oficial.

\section{El problema}
\label{sec:intro-problema}

El seguimiento multiobjeto consiste en asignar, a cada objeto detectado en un
vídeo, una identidad que se mantenga consistente a lo largo de todos los
frames en los que ese objeto sea visible, incluso cuando desaparece
temporalmente de la imagen, por ejemplo, al quedar oculto detrás de otro
jugador y vuelve a aparecer poco después. Un buen algoritmo de tracking debe
resolver simultáneamente dos problemas relacionados pero distintos:

\begin{itemize}
  \item \textbf{Detección}: encontrar dónde está cada objeto de interés en
        cada frame. En este proyecto, jugadores, árbitros y balón.
  \item \textbf{Asociación}: decidir qué detección de un frame corresponde a
        qué identidad ya existente de un frame anterior, o si se trata de un
        objeto nuevo.
\end{itemize}

Es importante delimitar aquí, desde el principio, el alcance exacto de este
proyecto: el dataset SoccerNet-Tracking proporciona las detecciones ya
calculadas para cada frame (\texttt{det.txt}). Por tanto, la tarea abordada es
específicamente la de \textbf{asociación}: mantener identidades
consistentes a partir de un conjunto de detecciones dado y no la de
detección de jugadores desde cero a partir de la imagen. Esta
decisión, natural dado el planteamiento del dataset, permite centrar el
esfuerzo y la comparativa en los algoritmos de tracking propiamente dichos
(el componente de asociación), que es donde reside la principal diferencia
entre los distintos métodos evaluados. Esta distinción se retoma con más
detalle en el capítulo~\ref{cap:analisis}.

\section{Objetivos}
\label{sec:intro-objetivos}

\subsection{Objetivo general}

Una vez definidas las detecciones del dataset SoccerNet, se propone evaluar y comparar distintos algoritmos de seguimiento multiobjeto sobre
secuencias reales de partidos de fútbol del dataset SoccerNet, y explorar
en qué medida las trayectorias resultantes permiten extraer metadatos
deportivos de utilidad para el análisis de un partido.

\subsection{Objetivos específicos}

A partir de este objetivo general se han definido los siguientes objetivos
específicos:

\begin{itemize}
  \item \textbf{OE-01:} Estudiar la estructura del dataset
        SoccerNet-Tracking (secuencias de frames, detecciones precomputadas
        y \emph{ground truth}) y delimitar la tarea abordada como un
        problema de asociación en mayor medida que de detección.
  \item \textbf{OE-02:} Implementar varios algoritmos de tracking de tipo
        \emph{tracking-by-detection}, incluyendo un
        filtrado geométrico que separe el balón de los jugadores antes de
        la asociación.
  \item \textbf{OE-03:} Ejecutar los algoritmos escogidos (\textbf{ByteTrack} y
        \textbf{OC-SORT}) sobre la totalidad de las
        secuencias de los conjuntos de entrenamiento y de test del dataset.
  \item \textbf{OE-04:} Evaluar ambos algoritmos con las métricas estándar
        del campo (HOTA, MOTA, IDF1) mediante \texttt{TrackEval}, comparando
        su comportamiento entre el conjunto de entrenamiento y el de test.
  \item \textbf{OE-05:} Realizar un estudio sistemático de hiperparámetros
        de cada algoritmo, variando un parámetro a la vez, y determinar una
        configuración óptima para cada uno de ellos.
  \item \textbf{OE-06:} Diseñar y aplicar un método de clasificación de
        jugadores por equipo a partir del color dominante de la camiseta
        mediante K-means.
  \item \textbf{OE-07:} Generar mapas de calor de ocupación del campo a
        partir de las trayectorias obtenidas, tanto de forma global como
        separados por equipo.
  \item \textbf{OE-08:} Calcular una homografía real sobre tramos de vídeo
        con cámara estable, como prueba de concepto para estimar la
        distancia recorrida y la velocidad de los jugadores en unidades
        reales.
  \item \textbf{OE-09:} Estimar la posesión del balón por proximidad
        espacial entre el balón y el jugador más cercano, y si es posible, 
        generar un resumen de posesión de balón por jugador y
        agregarla por equipo.
  \item \textbf{OE-10:} Documentar las limitaciones
        técnicas encontradas en cada etapa del proyecto, incluyendo
        enfoques explorados y descartados, de modo que la memoria refleje
        fielmente tanto lo conseguido como las dificultades reales del
        proceso.
  \item \textbf{OE-11:} Desarrollar una aplicación web que sirva como
        demostrador visual e interactivo de los resultados del proyecto obtenidos,
        permitiendo explorar los
        resultados de tracking y el análisis de datos generado. También se
        destinada a personas que no sean expertas en el área para facilitar
        su comprensión y visualización.
\end{itemize}

\section{Alcance y delimitación del proyecto}
\label{sec:intro-alcance}

Como se ha introducido en la sección~\ref{sec:intro-problema}, este proyecto
parte de las detecciones del dataset y se centra en la etapa de
asociación (tracking), no en el diseño de un detector de objetos propio. Estos son
 los límites principales del proyecto:

\begin{itemize}
  \item Los algoritmos comparados corresponden a la familia de métodos
        \emph{tracking-by-detection} clásicos, los usados no tienen componente de
        reidentificación visual profunda (\emph{re-ID}) pero la incorporación de
        algoritmos con re-ID se plantea como posible ampliación
        (capítulo~\ref{cap:conclusiones}).
  \item La proyección de posiciones de píxel a coordenadas reales del campo
        se realiza mediante dos enfoques de distinta fiabilidad: una
        aproximación por rango observado, usada
        de forma general, y una homografía real calculada manualmente sobre
        tramos concretos con cámara estable, usada como prueba de concepto
        para el cálculo de distancia y velocidad. Ambos enfoques, y sus
        limitaciones, se detallan en el capítulo~\ref{cap:implementacion}.
  \item La posesión del balón se estima por proximidad espacial entre el
        balón detectado y el jugador más cercano, una aproximación razonable
        pero distinta del concepto de posesión que manejaría un analista
        táctico humano.
\end{itemize}

\section{Estructura del documento}
\label{sec:intro-estructura}

La memoria de este proyecto se organiza en los siguientes capítulos:

\begin{itemize}
  \item \textbf{Capítulo~\ref{cap:fundamentos} -- Fundamentos teóricos y
        estado del arte.}
        Describe los conceptos teóricos necesarios para entender el
        proyecto y el estado del arte. Explica el problema del seguimiento multiobjeto, los algoritmos
        de tracking comparados, las métricas de evaluación empleadas, la
        clasificación de jugadores por color y los fundamentos de la
        homografía entre otros. y por último, cierra situando este TFG frente a los
        resultados del reto oficial de SoccerNet.
  \item \textbf{Capítulo~\ref{cap:gestion} -- Gestión del proyecto.}
        Describe la metodología de trabajo seguida, la planificación
        temporal, la estructura de descomposición del trabajo, y una
        estimación de los recursos y costes que habría requerido este
        proyecto en un contexto profesional.
  \item \textbf{Capítulo~\ref{cap:analisis} -- Análisis de requisitos.}
        Recoge los actores del sistema, los requisitos de información,
        funcionales y no funcionales, las reglas de negocio y los casos de
        uso.
  \item \textbf{Capítulo~\ref{cap:diseno} -- Diseño.}
        Presenta la arquitectura general del sistema y las decisiones de
        diseño tomadas.
  \item \textbf{Capítulo~\ref{cap:implementacion} -- Implementación.}
        Detalla cómo se han implementado el tracking, la evaluación, el
        estudio de hiperparámetros, el pipeline de metadatos y la
        aplicación web, tambien los problemas técnicos encontrados y los
        enfoques que se probaron y no funcionaron.
  \item \textbf{Capítulo~\ref{cap:pruebas} -- Resultados.}
        Recoge la comparativa final entre ByteTrack y OC-SORT, los
        resultados del estudio de hiperparámetros y del pipeline de
        metadatos, y las limitaciones detectadas en cada uno.
  \item \textbf{Capítulo~\ref{cap:manual} -- Manual de despliegue.}
        Explica paso a paso cómo instalar, configurar y ejecutar el
        sistema, incluida la aplicación web.
  \item \textbf{Capítulo~\ref{cap:uso-ia} -- Uso de inteligencia
        artificial.} Detalla el uso realizado de herramientas de
        inteligencia artificial durante el desarrollo del proyecto.
  \item \textbf{Capítulo~\ref{cap:conclusiones} -- Conclusiones y trabajo
        futuro.} Resume las
        conclusiones del proyecto y plantea posibles líneas futuras de
        trabajo.
\end{itemize}

Tras la bibliografía se incluye un glosario
con los términos y las siglas usadas a lo largo de la memoria.
